% Problem-section file following ../_template.tex.
\section{Quantum capacity of a bosonic thermal attenuator}
\label{sec:problem-49}

\paragraph{Problem.}
Let $\Phi_{\eta,\nu}$ be the single-mode bosonic thermal attenuator with
transmissivity $0<\eta<1$ and environmental mean photon number
$0<\nu<\infty$, defined by
\begin{equation}
  \Phi_{\eta,\nu}(\rho_A)
  :=\operatorname{Tr}_{E'}\!\left[
     U_\eta(\rho_A\otimes\tau_{\nu,E})U_\eta^\dagger
  \right],
  \qquad
  \tau_\nu:=\sum_{k=0}^{\infty}
     \frac{\nu^k}{(\nu+1)^{k+1}}|k\rangle\!\langle k|,
  \label{eq:p49-thermal-attenuator}
\end{equation}
where $U_\eta:AE\to BE'$ is a beam-splitter unitary.  Its unassisted
quantum capacity is the regularized coherent information
\begin{equation}
  \mathcal Q(\Phi_{\eta,\nu})
  :=\lim_{n\to\infty}\frac1n\sup_{\rho_{A^n}}
       I_{\rm c}(\rho_{A^n},\Phi_{\eta,\nu}^{\otimes n}),
  \qquad
  I_{\rm c}(\rho,\mathcal N)
  :=S(\mathcal N(\rho))-S(\mathcal N^{\rm c}(\rho)),
  \label{eq:p49-quantum-capacity}
\end{equation}
where $S(\sigma):=-\operatorname{Tr}(\sigma\log_2\sigma)$ and
$\mathcal N^{\rm c}$ is any complementary channel.  Determine
Eq.~\eqref{eq:p49-quantum-capacity} for all finite-temperature parameters in
Eq.~\eqref{eq:p49-thermal-attenuator}.

\subsection*{Status}

Unsolved

\subsection*{Source}

Rosati, Mari, and Giovannetti explicitly study the unknown thermal-attenuator
quantum capacity through nonmatching narrow bounds; the exact-capacity
problem is retained after the later non-Gaussian separation
\sourcecite{ref:p49-rosati-mari-giovannetti}{RMG18},
\sourcecite{ref:p49-mele-et-al}{MCF+26}.

\subsection*{Progress}

\begin{itemize}
  \item At zero environmental temperature, the pure-loss capacity is
  \begin{equation}
    \mathcal Q(\Phi_{\eta,0})
    =\left[\log_2\frac{\eta}{1-\eta}\right]_+,
    \qquad [x]_+:=\max\{0,x\}.
    \label{eq:p49-pure-loss-capacity}
  \end{equation}
  Thus Eq.~\eqref{eq:p49-pure-loss-capacity} solves only the boundary case
  $\nu=0$ \sourcecite{ref:p49-wolf-perez-garcia-giedke}{WPG07}.

  \item Optimizing the one-use coherent information over all single-mode
  Gaussian input states gives exactly
  \begin{equation}
    \sup_{\rho\in\mathsf G_1}I_{\rm c}(\rho,\Phi_{\eta,\nu})
    =\left[\log_2\frac{\eta}{1-\eta}-g(\nu)\right]_+,
    \qquad
    g(x):=(x+1)\log_2(x+1)-x\log_2x.
    \label{eq:p49-gaussian-coherent-information}
  \end{equation}
  The rate in Eq.~\eqref{eq:p49-gaussian-coherent-information} was first
  obtained from thermal inputs, but it is not the optimum over arbitrary
  non-Gaussian inputs
  \sourcecite{ref:p49-holevo-werner}{HW01},
  \sourcecite{ref:p49-mele-et-al}{MCF+26}.

  \item The exact antidegradability region is
  \begin{equation}
    \eta\leq\eta_{\rm AD}(\nu):=
    \frac{\nu+\frac12}{\nu+1},
    \qquad \mathcal Q(\Phi_{\eta,\nu})=0.
    \label{eq:p49-antidegradable-region}
  \end{equation}
  Hence Eq.~\eqref{eq:p49-quantum-capacity} is settled throughout the region
  in Eq.~\eqref{eq:p49-antidegradable-region}
  \sourcecite{ref:p49-lami-et-al}{LKA+19}.

  \item A bottleneck decomposition supplies a generally nonmatching upper
  bound:
  \begin{equation}
    \mathcal Q(\Phi_{\eta,\nu})
    \leq\left[
      \log_2\frac{\eta-(1-\eta)\nu}{(1-\eta)(\nu+1)}
    \right]_+
    \quad\text{when }\eta>(1-\eta)\nu,
    \label{eq:p49-bottleneck-bound}
  \end{equation}
  while the channel is entanglement breaking and has zero capacity otherwise
  \sourcecite{ref:p49-rosati-mari-giovannetti}{RMG18}.  The bound in
  Eq.~\eqref{eq:p49-bottleneck-bound} does not generally coincide with the
  Gaussian achievable rate.

  \item Non-Gaussian inputs strictly exceed the Gaussian value.  In
  particular,
  \begin{equation}
    \sup_{\rho\in\mathsf G_1}I_{\rm c}(\rho,\Phi_{4/5,1})=0,
    \qquad
    \mathcal Q(\Phi_{4/5,1})\geq4.7\times10^{-4}.
    \label{eq:p49-nongaussian-separation}
  \end{equation}
  Equation~\eqref{eq:p49-nongaussian-separation} rules out a restriction to
  single-mode Gaussian inputs but does not determine the capacity
  \sourcecite{ref:p49-mele-et-al}{MCF+26}.
\end{itemize}

\subsection*{References}

\begin{enumerate}[label={},leftmargin=4.8em,labelsep=0.5em]
  \footnotesize
  \item[\textup{[WPG07]}]\label{ref:p49-wolf-perez-garcia-giedke}
  M. M. Wolf, D. P\'erez-Garc\'ia, and G. Giedke,
  ``Quantum Capacities of Bosonic Channels,''
  \emph{Physical Review Letters} \textbf{98}, 130501 (2007).
  \href{https://doi.org/10.1103/PhysRevLett.98.130501}%
       {doi:10.1103/PhysRevLett.98.130501};
  \href{https://arxiv.org/abs/quant-ph/0606132}{arXiv:quant-ph/0606132}.

  \item[\textup{[HW01]}]\label{ref:p49-holevo-werner}
  A. S. Holevo and R. F. Werner,
  ``Evaluating Capacities of Bosonic Gaussian Channels,''
  \emph{Physical Review A} \textbf{63}, 032312 (2001).
  \href{https://doi.org/10.1103/PhysRevA.63.032312}%
       {doi:10.1103/PhysRevA.63.032312};
  \href{https://arxiv.org/abs/quant-ph/9912067}{arXiv:quant-ph/9912067}.

  \item[\textup{[MCF+26]}]\label{ref:p49-mele-et-al}
  F. A. Mele, G. Catalano, M. Fanizza, V. Giovannetti, and L. Lami,
  ``Bosonic Quantum Communication Beyond the Thermal Threshold,''
  arXiv preprint (2026).
  \href{https://arxiv.org/abs/2607.27449}{arXiv:2607.27449}.

  \item[\textup{[LKA+19]}]\label{ref:p49-lami-et-al}
  L. Lami, S. Khatri, G. Adesso, and M. M. Wilde,
  ``Extendibility of Bosonic Gaussian States,''
  \emph{Physical Review Letters} \textbf{123}, 050501 (2019).
  \href{https://doi.org/10.1103/PhysRevLett.123.050501}%
       {doi:10.1103/PhysRevLett.123.050501};
  \href{https://arxiv.org/abs/1904.02692}{arXiv:1904.02692}.

  \item[\textup{[RMG18]}]\label{ref:p49-rosati-mari-giovannetti}
  M. Rosati, A. Mari, and V. Giovannetti,
  ``Narrow Bounds for the Quantum Capacity of Thermal Attenuators,''
  \emph{Nature Communications} \textbf{9}, 4339 (2018).
  \href{https://doi.org/10.1038/s41467-018-06848-0}%
       {doi:10.1038/s41467-018-06848-0};
  \href{https://arxiv.org/abs/1801.04731}{arXiv:1801.04731}.
\end{enumerate}

\subsection*{Comment}

For generic $\nu>0$ outside the zero-capacity regions, the lower and upper
bounds do not coincide.  Non-Gaussian one-mode optimization and possible
multimode superadditivity both remain relevant to
Eq.~\eqref{eq:p49-quantum-capacity}.

\subsection*{Tag}

Bosonic channels; Quantum capacity; Coherent information;
Gaussian quantum information; Continuous-variable systems

\subsection*{ID}

\texttt{op\_9209e4eb15586dec}
